\documentclass[10pt,a4paper]{article}
\usepackage[utf8]{inputenc}
\usepackage[T1]{fontenc}
\usepackage{amsmath}
\usepackage{amssymb}
\usepackage{graphicx}
\usepackage{hyperref}
\title{Gattacca: An "unrealistic" Vision on the Future and Eugenics}
\author{R\^{i}mboi Eusebiu - grupa 332}
\begin{document}
	\maketitle
	
	I want to start my review on the movie by rating it personally at 6 out of 10, and there is a plenty of reasons behind this critique I will enumerate in a moment. But, for now, I will continue by presenting the good parts and let the bad ones to the end.
	
	By far, the most powerful idea of the film is that motivation, together with discipline (here manifested more like an obsession), beats, biologically speaking, even the best ones. Watching this moving and understanding the core idea it wants to send to the viewers, I can't agree more about it. Not only the discrimination based on genetics is very detrimental for the individual - who is constraint to leaving his/her aspirations and living how the norms allow depending on the genome quality, but also for the society as a whole, since a lot of inventors and revolutionaries we had, that changed entire nations and added huge contribution for the world we are lucky to live in right now, were not perfect, genetically speaking, sometimes not even close; but somehow they shown us that it's possible if you work hard and smart enough.
	
	Besides this, the other positive aspect I noticed in the movie was the genome identification of the employees. Even though the system had its problems and on some points we lack the technology to make that possible on such speed, its implementation was interesting to visualize and think of what parts of it could be implemented in the real life to increase the security of our systems and improve the efficiency of criminal investigations.
	
	Now, let's dive into the bad aspects. First, I will enunciate the core problems then enumerate some examples which showcase how these problem were observed by me.
	
	The first big issue with the film is, of course, the lack of realism. How is it possible for someone to fake such a system for so long? We don't talk about an ID card that was stole and the looks changed, no! We talk about an entire human genome that is faked! We see lots of countries and companies even nowadays that implement dozens of security measures that can detect such a fraud. Then, it comes the loss of Vincent's traces: how is it possible for a simple citizen, in such a sophisticated society to be able to disappear so easily from the government's sight? What's more questionable is how did Jerome escaped without any problems after beating to the ground a simple man that was doing its job. Lastly, there is the office atmosphere which intends to present the idea of perfection, tidiness, predictability and, at the same time, boredom. But it's a pretty wrong vision: these perfect people won't work for all the money in the world in such an office, under such dull working conditions. They know for sure a good workplace involves social activities, discussions, breaks and much more human stuff; not just some robotic environment and routines where the only "fun" activity is some kind of elitist opera. Besides this, working on such important projects, its crucial for everybody to work as a team since every small decision require multiple minds and very clear vision. Even though the people there are perfect, they can still make mistakes, some insignificant, others lethal; thus, it's in the benefit of literally everyone to work in an open space where the teamwork and communication is promoted. There are also a lot to say about how the investigations were held and the entrance interview, but it's enough for this point.
	
	The second point I want to make here is the crime scene. I understand they needed some action in the movie to add more anxiety and adrenaline in the watcher's mind, but the way they did that here was so unnatural and lacked of meaning that it looked ridiculous from my side: what do you mean you \textbf{killed} someone because that guy didn't want to approve you little space mission??? And how is it possible that those perfect detectives managed to catch the killer, not by using some lie detectors or talking to people or noticing any suspicious activity, but from a little stupid spittle??? After looking for every human being living on a radius of 10 km just for a guy that disappeared completely from the government's sight??? And what do you mean that detective got convinced by the real Jerome that he's in a wheelchair because some injury on a leg... 5 days before going into space??? Pathetic...
	
	There are also some considerations regarding the Vincent's heart condition that could kill him in space or even until reaching the orbit, but that's enough for this review.
	
	As a conclusion, the idea was inspiring, but how this movie was directed gives me hope that I can become myself a film director one day.

\end{document}